
The approach proceeds in two phases: (1) compression response profiling to generate a task-configuration matrix, and (2) statistical analysis to determine whether tasks cluster and whether attention features predict cluster membership.

\subsubsection{3.1 Compression Response Profiling}

Six compression configurations span eviction and quantization dimensions:

\begin{table}[htbp]
\centering
\resizebox{\columnwidth}{!}{%
\begin{tabular}{llll}
\toprule
Config & Method & Retention & Quantization \\
\midrule
C1 & Full & 100\% & FP16 \\
C2 & H2O & 80\% & FP16 \\
C3 & H2O & 40\% & FP16 \\
C4 & Full & 100\% & INT8 \\
C5 & Full & 100\% & INT4 \\
C6 & H2O & 60\% & INT8 \\
\bottomrule
\end{tabular}
}
\end{table}

For each of 21 LongBench tasks t and each configuration c, accuracy A(t, c) is computed using task-appropriate metrics. The response matrix R $\in \mathbb{R}^{21}ˣ^6$ contains accuracy retention ratios:

R\_tc = A(t, c) / A(t, c\_full)

\subsubsection{3.2 Cluster Discovery via Gap Statistic}

The gap statistic determines the optimal number of clusters k* directly from data, avoiding bias from assuming task-compression structure exists.

For k = 1, ..., K\_max:
\begin{enumerate}
\item Cluster R into k groups using k-means, computing within-cluster dispersion $W_k$
\item Generate B reference datasets from uniform distribution on the same bounding box
\item Compute Gap(k) = E*[log $W_k$] - log $W_k$
\end{enumerate}

The optimal k* is the smallest k satisfying: Gap(k) >= Gap(k+1) - $s_{k+1}$

Parameters: B = 500 bootstrap samples, K\_max = 6.

\subsubsection{3.3 Attention Entropy Analysis}

For each task, attention entropy is computed from the first 100 tokens across all 32 layers and 32 heads:

$H(\alpha ) = -\Sigma _i \alpha _i log(\alpha _i + \epsilon )$

where $\epsilon = 10^{-10}$. One-way ANOVA tests whether entropy differs across 6 LongBench task domains with entropy as the dependent variable.

\subsubsection{3.4 Entropy-Tolerance Relationship}

A two-sample t-test compares accuracy retention under 40\% H2O eviction between high-entropy domains (Single-Document QA, Multi-Document QA, Long-dialogue History Understanding) and low-entropy domains (Long Structured Data Understanding, Long In-context Learning, Code Repository Understanding), based on median split of domain mean entropy values.
